\chapterimage{figs/chapterhead.png}
\chapter{Model Editing in the GUI}
\label{chap:model-editing-in-the-gui}

This chapter turns the GUI from a passive viewer into an editing workspace. The goal is to show a
reproducible path for creating or loading a model, changing it, checking it, saving it, and opening
it again without losing the structure that was just edited.

The safest way to follow the chapter is to start from a checked-in example model such as
\path{models/AirportSecurityExample.gen}, save a working copy, and make the edits in that copy. That
keeps the original file intact and makes the round-trip easier to verify.

The chapter is organized around these steps:
Figure~\ref{fig:model-editing-workflow},
Figure~\ref{fig:model-editing-connection},
Figure~\ref{fig:model-editing-invalid-property},
Figure~\ref{fig:model-editing-validation-message}, and
Figure~\ref{fig:model-editing-save-reopen}.

\noindent\textbf{Objective.} Show how a user can create, connect, edit, validate, save, and reopen a
model through the graphical interface.

\noindent\textbf{Audience.} Readers who already know the main GUI layout and now need to perform real
editing work on a model.

\noindent\textbf{Scope.} Cover the component insertion path, connection flow, property editing,
duplicate/remove behavior, validation, and the persistence round-trip back to disk.

The implementation used by the GUI is intentionally conservative. The current code exposes model
creation, open, save, check, undo, redo, cut, copy, paste, delete, group, ungroup, and replace
actions. There is no separate dedicated ``duplicate'' command in the current GUI path, so this chapter
treats duplication as the documented copy-paste workflow.

\section{Starting from a reproducible baseline}
\label{sec:model-editing-baseline}

Editing should begin from a stable model state. In practice, that means one of two things:
\begin{itemize}
\item create a new model from the GUI template;
\item open an existing checked-in model and save it under a new name before editing.
\end{itemize}

For the purposes of this chapter, the second option is the better baseline. A checked-in example gives
the reader a known starting point, a known structure, and a known file to compare before and after the
edits.

The relevant GUI entry points are the model lifecycle actions wired in the main window:
\begin{itemize}
\item \texttt{Model New};
\item \texttt{Model Open};
\item \texttt{Model Save};
\item \texttt{Model Check}.
\end{itemize}

The same workspace also exposes the property editor, the component tree, the graphical scene, and the
current model tabs. That combination is what makes editing possible without leaving the application.

\begin{figure}[htbp]
\centering
% FIGURE-SPEC-BEGIN
% ID: fig:model-editing-workflow
% Source of truth:
%   - models/AirportSecurityExample.gen
%   - source/applications/gui/genesys/mainwindow.ui
%   - source/applications/gui/genesys/mainwindow_controller.cpp
%   - source/applications/gui/genesys/controllers/ModelLifecycleController.cpp
%   - source/applications/gui/genesys/controllers/PropertyEditorController.cpp
%   - source/applications/gui/genesys/controllers/EditCommandController.cpp
%   - source/applications/gui/genesys/mainwindow.cpp
% Purpose:
%   Show the practical editing loop from loading a model to selecting elements and changing properties.
% Required visual content:
%   - the loaded model canvas
%   - the component tree or palette used to reach components
%   - the property editor bound to a selected component
%   - a visible indication that the model is in an editable state
% Accessibility description:
%   A workspace screenshot that proves the GUI is ready for model editing.
% Validation:
%   The screen must correspond to the checked-in model or its saved working copy and not to an unrelated scene.
% FIGURE-SPEC-END
\figureplaceholder{figs/generated/model-editing-workflow}
\caption{Workspace ready for model editing.}
\label{fig:model-editing-workflow}
\end{figure}

Figure~\ref{fig:model-editing-workflow} is the anchor for the rest of the chapter. Every later edit,
connection, validation, and persistence step should be read against the same loaded model.

\section{Creating components and placing them in the model}
\label{sec:model-editing-creating-components}

Once a model is open, editing begins by inserting or selecting components in the graphical scene. The
GUI keeps the scene, the property editor, and the tree views synchronized so that selecting an element
immediately exposes its editable fields.

The practical sequence is simple:
\begin{enumerate}
\item choose the model or component source from the GUI;
\item place a component in the scene;
\item select the new component;
\item confirm that the property editor follows the selection;
\item repeat until the desired flow exists.
\end{enumerate}

If the user starts from a blank model, the same cycle applies, only with more insertions. If the user
starts from the checked-in tutorial model, the cycle becomes a controlled edit to an existing flow.

The main point is that creation is not separate from inspection. A component should be created and then
immediately checked in the property editor, because most modeling errors start as a mismatch between the
intended role of the component and its configured properties.

\begin{figure}[htbp]
\centering
% FIGURE-SPEC-BEGIN
% ID: fig:model-editing-connection
% Source of truth:
%   - source/applications/gui/genesys/graphicals/ModelGraphicsScene.h
%   - source/applications/gui/genesys/graphicals/ModelGraphicsView.h
%   - source/applications/gui/genesys/controllers/EditCommandController.cpp
%   - source/applications/gui/genesys/mainwindow.ui
% Purpose:
%   Show how one model component is connected to another in the graphical scene.
% Required visual content:
%   - two or more components with a visible connection path
%   - source and destination ports or connection handles
%   - the current scene state after the connection is made
% Accessibility description:
%   A close-up that makes the direction and ownership of the connection easy to understand.
% Validation:
%   The connection must be a real scene link between live components, not just a diagram sketch.
% FIGURE-SPEC-END
\figureplaceholder{figs/generated/model-editing-connection}
\caption{Connection between model components.}
\label{fig:model-editing-connection}
\end{figure}

Figure~\ref{fig:model-editing-connection} should show the connection workflow as it exists in the scene,
not as an abstract diagram. The user should be able to recognize source and destination elements and
see that the connection belongs to the current model.

\section{Editing properties and provoking one controlled error}
\label{sec:model-editing-properties-and-validation}

Property editing is where the model becomes concrete. In the GUI, selecting a component or data
definition updates the property editor so the user can change names, numeric parameters, expressions,
and other configuration fields.

The recommended sequence is:
\begin{enumerate}
\item select one component;
\item change one safe property first;
\item deliberately introduce one invalid property value;
\item run \texttt{Model Check};
\item read the error feedback in the console or dialog;
\item correct the value and check again.
\end{enumerate}

This is not a decorative exercise. The point is to prove that the GUI does more than store the visual
layout: it also validates the model before execution. The code path behind \texttt{Model Check} calls the
model's own \texttt{check()} logic, refreshes the graphical data definitions when needed, and marks the
model as valid or invalid based on the result.

The chapter does not invent a separate validation mechanism. It uses the same check action that the GUI
already exposes.

\begin{figure}[htbp]
\centering
% FIGURE-SPEC-BEGIN
% ID: fig:model-editing-invalid-property
% Source of truth:
%   - source/applications/gui/genesys/controllers/PropertyEditorController.cpp
%   - source/applications/gui/genesys/mainwindow.cpp
%   - source/applications/gui/genesys/mainwindow.ui
%   - models/AirportSecurityExample.gen
% Purpose:
%   Show one deliberately invalid property value before validation is run.
% Required visual content:
%   - the selected component in the property editor
%   - one property field containing an invalid or incomplete value
%   - the edited scene in the background or nearby
% Accessibility description:
%   The invalid state must be visually obvious without needing to inspect raw files.
% Validation:
%   The property shown must be one that the model actually exposes in the current GUI.
% FIGURE-SPEC-END
\figureplaceholder{figs/generated/model-editing-invalid-property}
\caption{A deliberately invalid property value before validation.}
\label{fig:model-editing-invalid-property}
\end{figure}

\begin{figure}[htbp]
\centering
% FIGURE-SPEC-BEGIN
% ID: fig:model-editing-validation-message
% Source of truth:
%   - source/applications/gui/genesys/mainwindow.cpp
%   - source/applications/gui/genesys/controllers/ModelLifecycleController.cpp
%   - source/kernel/simulator/model/Model.cpp
% Purpose:
%   Show the model validation result after the invalid property is checked.
% Required visual content:
%   - the validation dialog or console message
%   - the model check action in context
%   - the error text or success/failure indication
% Accessibility description:
%   A visual proof that the GUI reports validation results to the user.
% Validation:
%   The message must correspond to the model's current check result, not a previous run.
% FIGURE-SPEC-END
\figureplaceholder{figs/generated/model-editing-validation-message}
\caption{Model validation feedback after a check.}
\label{fig:model-editing-validation-message}
\end{figure}

Figure~\ref{fig:model-editing-invalid-property} shows the bad state. Figure~\ref{fig:model-editing-validation-message}
shows the result of asking the application to validate it. Together they prove that the editor and the
checker are coupled in the expected way.

\section{Undo, redo, copy, paste, delete, group, and ungroup}
\label{sec:model-editing-edit-operations}

Once a model is editable, the GUI needs to support routine editing operations. The current code exposes
them through the edit menu and the scene commands:
\begin{itemize}
\item undo;
\item redo;
\item cut;
\item copy;
\item paste;
\item delete;
\item group;
\item ungroup;
\item replace.
\end{itemize}

The most important practical point is that copy and paste are the current conservative path for
duplication. There is no separate dedicated duplicate action in the code path inspected for this chapter.
If a reader wants to duplicate a model element, the supported route is to copy it and then paste it into
the same model or a working copy.

Deletion is equally direct. The scene filters the selected items to the ones that can actually be
manipulated, and the delete command removes only those items. Grouping and ungrouping are available for
layout management, but they do not replace the underlying model relationships.

These commands matter because editing a model is usually a sequence of small corrections rather than one
big creation step. A user adds, repositions, duplicates, and removes elements until the flow matches the
intended logic.

\section{Saving and reopening the edited model}
\label{sec:model-editing-save-and-reopen}

After the model is in the desired state, save it and reopen it immediately. That is the fastest way to
prove the persistence round-trip.

The GUI accepts a graphical save target with the \texttt{.gui} extension and also supports saving a text
model to \texttt{.gen}. For an editing workflow, \texttt{.gui} is the conservative choice because it
preserves the graphical state that was just changed. If the user needs a kernel-oriented artifact, the
\texttt{.gen} route remains available.

The practical verification steps are:
\begin{enumerate}
\item save the working copy under a new base name;
\item close or switch away from the model;
\item reopen the saved file;
\item confirm that the scene and properties match the saved state;
\item verify that the model is no longer marked as dirty.
\end{enumerate}

This is the point where the chapter proves persistence rather than just editing. If the model opens with
the same structure and the same properties, the GUI round-trip worked.

\begin{figure}[htbp]
\centering
% FIGURE-SPEC-BEGIN
% ID: fig:model-editing-save-reopen
% Source of truth:
%   - source/applications/gui/genesys/controllers/ModelLifecycleController.cpp
%   - source/applications/gui/genesys/mainwindow.cpp
%   - source/applications/gui/genesys/mainwindow_modelrepresentations.cpp
%   - models/AirportSecurityExample.gen
% Purpose:
%   Show the edited model after saving and after reopening the saved copy.
% Required visual content:
%   - the saved file name or tab title
%   - the reopened scene
%   - the property editor or model tree confirming the restored state
%   - a visible indication that the round-trip succeeded
% Accessibility description:
%   A before-and-after proof that the edited model persisted correctly.
% Validation:
%   The reopened model must match the saved working copy, not the original baseline file.
% FIGURE-SPEC-END
\figureplaceholder{figs/generated/model-editing-save-reopen}
\caption{Edited model after save and reopen.}
\label{fig:model-editing-save-reopen}
\end{figure}

Figure~\ref{fig:model-editing-save-reopen} is the final proof that the editing session was not temporary.
If the saved file can be reopened and the change remains visible, the workflow is complete.

\section{What this chapter proves}
\label{sec:model-editing-what-it-proves}

This chapter establishes a narrow but useful claim:
\begin{itemize}
\item the GUI can create or open a model;
\item the model can be edited through the scene and property editor;
\item the GUI can expose and report validation problems;
\item changes can be saved and reopened;
\item the same working copy survives the persistence round-trip.
\end{itemize}

The chapter does not claim that every optional export path, advanced plugin, or experimental workflow is
fully documented here. It focuses on the editing path that is already visible in the current GUI and on
the behavior that can be confirmed directly from the repository.

That is enough to move to the next chapters, where the manual starts explaining execution, results, and
the more specialized model workflows that build on this editing foundation.

Test your understanding of this content by answering the following questions:

\begin{enumerate}
\item Why is it safer to start from a checked-in model and save a working copy before editing?

\item Which GUI actions form the core edit loop in this chapter?

\item Why does the chapter treat duplication as copy and paste instead of a separate command?

\item What proves that the edited model was really saved and reopened?
\end{enumerate}
